\chapter{Résultats, discussion et limites de l'étude}
\label{chap:resultats}

Ce chapitre confronte le cadre méthodologique du Chapitre~\ref{chap:methodo}
aux données réellement collectées. Nous y lisons les exports du pipeline sous
l'angle des questions $Q_1$--$Q_3$ et des hypothèses $H_1$--$H_3$ formulées en
introduction. L'enjeu n'est pas d'accumuler des tableaux, mais de montrer ce
que le protocole permet --- et ce qu'il interdit --- de conclure.

La campagne analytique porte sur la période du 15~juin au 15~juillet~2026
(UTC). Après sélection des stations outdoor actives et passage du contrôle
qualité, le corpus compte 38~stations et 2\,445\,899~observations de
température, d'humidité relative et de pression. Aucune référence
métrologique colocalisée n'étant disponible, nous n'affichons ni MAE, ni RMSE,
ni MAPE, ni $R^2$ comme preuves d'exactitude. Le refus est méthodologique :
ces indicateurs mesureraient un écart à une source non équivalente, non une
erreur de mesure démontrée.

\section{Fil de lecture du chapitre}
\label{sec:fil}

Le travail vise un cadre reproductible d'évaluation de la confiance pour des
données de stations personnelles low-cost. L'application empirique s'appuie sur
des stations openSenseMap qui exposent simultanément température, humidité et
pression.

Trois niveaux d'analyse sont maintenus distincts. Au niveau de l'observation,
nous examinons drapeaux et statut QC. Au niveau de la série (station et
variable), nous calculons dimensions et score de confiance. Au niveau de
l'usage, nous attribuons une classe \textit{fit for purpose}, puis nous
agrégeons cette classe à la station selon la règle du minimum.

Les résultats répondent à trois questions. ($Q_1$)~Le corpus est-il assez dense
et cohérent pour un diagnostic de confiance~? ($Q_2$)~Le découpage en dimensions
sépare-t-il des profils de séries~? ($Q_3$)~La classification d'usage
apporte-t-elle une information plus restrictive qu'un score global unique~?

Trois hypothèses orientent l'interprétation. ($H_1$)~Une part notable des
observations sera marquée suspecte, sans invalider l'ensemble du corpus.
($H_2$)~Les variables n'auront pas le même profil ; l'humidité sera plus
sensible aux problèmes de stabilité. ($H_3$)~Un score médian élevé pourra
coexister avec une classe d'usage modeste dès qu'une seule variable échoue.

%==============================================================================
\section{Synthèse du corpus}
\label{sec:corpus}

\subsection{Faits observés}

Le Tableau~\ref{tab:corpus} résume le corpus retenu après sélection des stations
et application du contrôle qualité.

\begin{table}[H]
\centering
\caption{Synthèse du corpus empirique (exports du pipeline).}
\label{tab:corpus}
\begin{tabular}{lc}
\toprule
Indicateur & Valeur \\
\midrule
Nombre de stations & 38 \\
Nombre de séries (station et variable) & 114 \\
Variables & température, humidité, pression \\
Observations après QC & 2\,445\,899 \\
Complétude médiane & 93,88\,\% \\
Score de confiance médian (séries) & 0,965 \\
Intervalle interquartile du score & 0,907 à 0,987 \\
Contrôle spatial & désactivé \\
\bottomrule
\end{tabular}
\end{table}

Les 114 séries correspondent à $38 \times 3$ couples station--variable. Avec une
complétude médiane proche de 94\,\%, la moitié des séries conserve, sur la
période, une couverture temporelle élevée une fois l'intervalle médian
d'échantillonnage pris comme référence. Le score de confiance, lui, n'est pas
symétrique : la médiane (0,965) dépasse la moyenne (0,923), l'écart-type vaut
0,100 et le minimum s'établit à 0,575 (Tableau~\ref{tab:trust_desc}). Autrement
dit, un noyau de séries très bien notées coexiste avec une traîne de séries
nettement plus faibles.

\begin{table}[H]
\centering
\caption{Distribution du score de confiance au niveau des séries ($n=114$).}
\label{tab:trust_desc}
\begin{tabular}{lcccccc}
\toprule
 & $n$ & Moyenne & Écart-type & Min & Médiane & Max \\
\midrule
\texttt{trust\_score} & 114 & 0,923 & 0,100 & 0,575 & 0,965 & 0,995 \\
\bottomrule
\end{tabular}
\end{table}

La Figure~\ref{fig:dist_trust} visualise cette concentration dans le haut de
l'échelle $[0,1]$, ainsi que la dispersion par variable.

\begin{figure}[H]
\centering
\includegraphics[width=0.95\textwidth]{fig1_distribution_trust_scores.png}
\caption{Distribution des scores de confiance au niveau des séries (gauche) et
par variable (droite).}
\label{fig:dist_trust}
\end{figure}

\subsection{Interprétation}

Ces ordres de grandeur ne disent pas que les données sont exactes au sens
métrologique. Ils disent que, \textit{sous les règles explicites du cadre}, la
majorité des séries combine disponibilité correcte, plausibilité physique élevée
et continuité temporelle satisfaisante. L'écart médiane--moyenne rappelle toutefois
qu'un sous-ensemble dégradé tire la moyenne vers le bas. Un seul indicateur
central ne suffit donc pas à décrire le corpus.

Du point de vue de $Q_1$, le matériau est suffisant pour exercer le protocole :
volume important, trois variables homogènes, couverture globalement dense. Reste
à localiser les dégradations.

%==============================================================================
\section{Contrôle qualité au niveau observation}
\label{sec:qc}

\subsection{Répartition des statuts \texttt{valide}, \texttt{a\_verifier}, \texttt{suspecte}}

Sur les 2\,445\,899 observations soumises au QC
(Tableau~\ref{tab:qc_status} et Figure~\ref{fig:qc_status}),
1\,859\,834~sont classées \texttt{valide} (76,0\,\%),
14\,307~\texttt{a\_verifier} (0,6\,\%) et 571\,758~\texttt{suspecte}
(23,4\,\%).

\begin{table}[H]
\centering
\caption{Statuts QC au niveau observation.}
\label{tab:qc_status}
\begin{tabular}{lrr}
\toprule
Statut & Observations & Proportion \\
\midrule
\texttt{valide} & 1\,859\,834 & 76,04\,\% \\
\texttt{a\_verifier} & 14\,307 & 0,58\,\% \\
\texttt{suspecte} & 571\,758 & 23,38\,\% \\
\bottomrule
\end{tabular}
\end{table}

\begin{figure}[H]
\centering
\includegraphics[width=0.95\textwidth]{fig2b_qc_status_global.png}
\caption{Répartition globale des statuts de contrôle qualité.}
\label{fig:qc_status}
\end{figure}

La classe \texttt{a\_verifier} est marginale. En revanche, près d'un quart des
observations porte au moins un drapeau majeur --- plausibilité, valeur figée,
saut temporel ou anomalie statistique robuste. $H_1$ trouve ici un appui net :
le cadre isole une masse non négligeable d'observations douteuses sans rejeter
l'ensemble du jeu. Un usage direct des brutes, sans filtrage, serait donc
difficilement défendable.

\subsection{Différences entre variables}

Le croisement variable $\times$ statut (Tableau~\ref{tab:qc_var} et
Figure~\ref{fig:qc_var}) nuance le tableau global. La température concentre
davantage d'observations \texttt{valide} (681\,836) que l'humidité (584\,175) ou
la pression (593\,823), pour des volumes totaux proches. Humidité et pression
regroupent chacune environ 227\,000 observations \texttt{suspecte}, contre
environ 117\,000 pour la température.

\begin{table}[H]
\centering
\caption{Effectifs des statuts QC par variable.}
\label{tab:qc_var}
\begin{tabular}{lrrr}
\toprule
Variable & \texttt{valide} & \texttt{a\_verifier} & \texttt{suspecte} \\
\midrule
Humidité & 584\,175 & 4\,324 & 226\,637 \\
Pression & 593\,823 & 4\,074 & 227\,957 \\
Température & 681\,836 & 5\,909 & 117\,164 \\
\bottomrule
\end{tabular}
\end{table}

\begin{figure}[H]
\centering
\includegraphics[width=0.95\textwidth]{fig2c_qc_status_par_variable.png}
\caption{Statuts QC par variable : proportions empilées et effectifs.}
\label{fig:qc_var}
\end{figure}

Les taux médians d'anomalies au niveau série (Tableau~\ref{tab:qc_rates})
identifient le mécanisme dominant. Pour l'humidité, la médiane du taux de
valeurs figées (\texttt{stuck}) atteint 0,126, contre 0 pour la température et
la pression. Les sauts temporels restent faibles (médianes voisines de 0,01).
Les anomalies physiques et statistiques médianes sont nulles. Le problème
principal n'est donc pas une explosion de valeurs hors plage, mais une
\textbf{stabilité dégradée} sur une partie des séries d'humidité.

\begin{table}[H]
\centering
\caption{Indicateurs QC médians par variable (niveau série).}
\label{tab:qc_rates}
\begin{tabular}{lcccccc}
\toprule
Variable & Complétude & Physique & Figé & Saut & Valide & Suspecte \\
\midrule
Humidité & 0,938 & 0,000 & 0,126 & 0,013 & 0,850 & 0,143 \\
Pression & 0,939 & 0,000 & 0,000 & 0,010 & 0,978 & 0,018 \\
Température & 0,939 & 0,000 & 0,000 & 0,013 & 0,973 & 0,018 \\
\bottomrule
\end{tabular}
\end{table}

\begin{figure}[H]
\centering
\includegraphics[width=0.92\textwidth]{fig7_taux_anomalies_par_variable.png}
\caption{Médianes des taux d'anomalies QC par variable.}
\label{fig:anom_var}
\end{figure}

\subsection{Lecture scientifique}

Le statut QC répond à une question précise : l'observation est-elle compatible
avec les règles de cohérence interne retenues~? Il ne répond pas à celle de
l'exactitude par rapport à une référence étalonnée. Une observation
\texttt{suspecte} est une observation à traiter avec prudence, non une erreur
prouvée.

L'asymétrie humidité / température--pression est compatible avec le
comportement souvent rapporté des capteurs low-cost de type BME ou HTU en
environnement réel --- dérives, saturation, mauvaise ventilation. Les données
ne permettent pourtant pas de trancher entre cause matérielle, installation ou
microclimat. Cette limite causale revient en Section~\ref{sec:limites}.

%==============================================================================
\section{Scores dimensionnels et score global de confiance}
\label{sec:confiance}

\subsection{Rappel de construction}

Pour chaque série, six dimensions dans $[0,1]$ sont dérivées des taux QC :
complétude, plausibilité physique, stabilité, anomalies statistiques,
continuité, métadonnées. Le score global est une moyenne pondérée renormalisée
sur les dimensions disponibles. Poids retenus : complétude~0,20 ; physique~0,25 ;
stabilité~0,20 ; statistique~0,15 ; continuité~0,10 ; métadonnées~0,05. La
dimension spatiale est absente, le contrôle de voisinage ayant été désactivé.

\subsection{Profils observés par variable}

Les médianes des dimensions (Tableau~\ref{tab:dims} et
Figure~\ref{fig:profil_dim}) dessinent un profil net. Plausibilité physique,
qualité statistique et métadonnées restent excellentes pour les trois variables,
avec des médianes proches de~1. La continuité est très élevée (environ 0,997) et
la complétude se situe autour de 0,94. La stabilité, en revanche, sépare
l'humidité (médiane 0,929) de la pression (0,993) et de la température (0,992).

Le score médian suit le même ordre : température (0,976), pression (0,969), puis
humidité (0,955).

\begin{table}[H]
\centering
\caption{Médianes des dimensions de confiance et du score global par variable.}
\label{tab:dims}
\begin{tabular}{lccccccc}
\toprule
Variable & Compl. & Phys. & Stab. & Stat. & Cont. & Métad. & Score \\
\midrule
Humidité & 0,938 & 1,000 & 0,929 & 1,000 & 0,997 & 1,000 & 0,955 \\
Pression & 0,939 & 1,000 & 0,993 & 1,000 & 0,997 & 1,000 & 0,969 \\
Température & 0,939 & 1,000 & 0,992 & 1,000 & 0,997 & 1,000 & 0,976 \\
\bottomrule
\end{tabular}
\end{table}

\begin{figure}[H]
\centering
\includegraphics[width=0.92\textwidth]{fig2_profil_dimensions.png}
\caption{Profil moyen des dimensions de confiance par variable.}
\label{fig:profil_dim}
\end{figure}

\subsection{Lien entre complétude et confiance}

La Figure~\ref{fig:comp_trust} confronte taux de complétude et score. Le
coefficient de Spearman vaut $\rho = 0,59$ ($n=114$). La relation est positive
et modérée : la disponibilité contribue au score, mais ne le détermine pas. Une
série relativement complète peut être pénalisée par la stabilité ; à l'inverse,
une complétude imparfaite n'interdit pas un score élevé si les autres dimensions
restent saines. Cette dissociation justifie, a posteriori, de ne pas réduire la
confiance à un simple taux de remplissage.

\begin{figure}[H]
\centering
\includegraphics[width=0.72\textwidth]{fig4_completude_vs_confiance.png}
\caption{Relation entre taux de complétude et score de confiance, par variable.}
\label{fig:comp_trust}
\end{figure}

\subsection{Retour sur $Q_2$ et $H_2$}

Le cadre dimensionnel remplit sa fonction différenciante. Il ne se contente pas
d'un taux global d'« anomalies » : il montre \textit{où} la confiance se
dégrade. Ici, la dégradation dominante concerne la stabilité de l'humidité, ce
qui confirme $H_2$. Pour un usage scientifique prudent, cette lecture est plus
utile qu'un indicateur unique opaque.

%==============================================================================
\section{Classification \textit{fit for purpose}}
\label{sec:ffp}

\subsection{Résultats au niveau des séries}

La classification d'usage attribue à chaque série la classe la plus exigeante
dont tous les seuils sont respectés ; à défaut, la série est jugée
\texttt{non\_recommande}. Sur 114~séries (Tableau~\ref{tab:ffp_series} et
Figure~\ref{fig:ffp_effectifs}), 55~(48,2\,\%) relèvent de
\texttt{usage\_scientifique\_restreint}, 16~(14,0\,\%) de
\texttt{comparaison\_relative}, 9~(7,9\,\%) de \texttt{suivi\_tendances},
21~(18,4\,\%) d'\texttt{exploratoire} et 13~(11,4\,\%) de
\texttt{non\_recommande}.

\begin{table}[H]
\centering
\caption{Répartition des classes \textit{fit for purpose} (séries).}
\label{tab:ffp_series}
\begin{tabular}{lrr}
\toprule
Classe & $n$ séries & Proportion \\
\midrule
Usage scientifique restreint & 55 & 48,2\,\% \\
Comparaison relative & 16 & 14,0\,\% \\
Suivi de tendances & 9 & 7,9\,\% \\
Exploratoire & 21 & 18,4\,\% \\
Non recommandé & 13 & 11,4\,\% \\
\bottomrule
\end{tabular}
\end{table}

\begin{figure}[H]
\centering
\includegraphics[width=0.95\textwidth]{fig3a_fit_for_purpose_effectifs.png}
\caption{Effectifs des classes \textit{fit for purpose} : séries et stations.}
\label{fig:ffp_effectifs}
\end{figure}

Près de la moitié des séries atteint la classe la plus exigeante du barème. Dans
le même temps, près de 30\,\% restent cantonnées aux usages exploratoires ou non
recommandés. Le cadre n'absout pas le corpus : il hiérarchise des usages
possibles.

\subsection{Différences par variable}

Le Tableau~\ref{tab:ffp_var} et la Figure~\ref{fig:ffp_var} montrent que la
température est la plus souvent classée en usage scientifique restreint
(21/38, soit 55\,\%), suivie de la pression (19/38, 50\,\%) puis de l'humidité
(15/38, 39\,\%). L'humidité alimente surtout la classe exploratoire (12~séries).
La pression concentre davantage de séries \texttt{non\_recommande} (8) que la
température (4) ou l'humidité (1). Les échecs d'usage pour la pression ne se
réduisent donc pas au seul phénomène des valeurs figées ; d'autres dimensions,
notamment physiques, y interviennent.

\begin{table}[H]
\centering
\caption{Classes FFP par variable (effectifs de séries).}
\label{tab:ffp_var}
\begin{tabular}{lccccc}
\toprule
Variable & Sci. restreint & Comp. relative & Tendances & Exploratoire & Non recomm. \\
\midrule
Humidité & 15 & 9 & 1 & 12 & 1 \\
Pression & 19 & 3 & 4 & 4 & 8 \\
Température & 21 & 4 & 4 & 5 & 4 \\
\bottomrule
\end{tabular}
\end{table}

\begin{figure}[H]
\centering
\includegraphics[width=0.95\textwidth]{fig3b_ffp_par_variable.png}
\caption{Répartition des classes FFP par variable.}
\label{fig:ffp_var}
\end{figure}

\subsection{Score versus classe}

La Figure~\ref{fig:score_classe} indique que les scores tendent à décroître
lorsque l'on descend dans la hiérarchie d'usage, sans chevauchement total. Une
série peut conserver un score encore élevé tout en échouant à un seuil dur ---
exigence de stabilité ou de plausibilité pour la classe scientifique, par
exemple. C'est précisément l'intérêt de la couche FFP par rapport au seul
\texttt{trust\_score}.

\begin{figure}[H]
\centering
\includegraphics[width=0.95\textwidth]{fig3c_score_par_classe_ffp.png}
\caption{Score de confiance selon la classe FFP (boxplot et dispersion).}
\label{fig:score_classe}
\end{figure}

Le profil dimensionnel moyen par classe (Figure~\ref{fig:profil_ffp} et
Tableau~\ref{tab:profil_ffp}) confirme le mécanisme. Les séries
\texttt{usage\_scientifique\_restreint} présentent des moyennes élevées sur
toutes les dimensions. Les séries \texttt{non\_recommande} se distinguent
surtout par un effondrement de la dimension physique moyenne (0,23). La classe
\texttt{exploratoire} est marquée par une stabilité moyenne dégradée (0,70).
Deux voies d'échec se dessinent ainsi : l'une plutôt physique, l'autre plutôt
liée à la stabilité.

\begin{table}[H]
\centering
\caption{Profil dimensionnel moyen par classe FFP.}
\label{tab:profil_ffp}
\begin{tabular}{lcccccc}
\toprule
Classe & Compl. & Phys. & Stab. & Stat. & Cont. & Métad. \\
\midrule
Usage sci. restreint & 0,950 & 0,999 & 0,988 & 1,000 & 0,996 & 1,000 \\
Comparaison relative & 0,833 & 0,996 & 0,948 & 1,000 & 0,992 & 1,000 \\
Suivi tendances & 0,736 & 0,965 & 0,983 & 0,962 & 0,993 & 1,000 \\
Exploratoire & 0,799 & 0,965 & 0,702 & 0,973 & 0,977 & 1,000 \\
Non recommandé & 0,741 & 0,231 & 0,776 & 0,979 & 0,960 & 1,000 \\
\bottomrule
\end{tabular}
\end{table}

\begin{figure}[H]
\centering
\includegraphics[width=0.95\textwidth]{fig6_profil_dimensions_par_classe.png}
\caption{Profil moyen des dimensions de confiance par classe FFP.}
\label{fig:profil_ffp}
\end{figure}

La Figure~\ref{fig:stab_phys} situe les séries dans le plan stabilité--plausibilité.
Les seuils associés à l'usage scientifique restreint (plausibilité
$\geq 0,97$, stabilité $\geq 0,90$) séparent une partie du nuage ; les séries
\texttt{non\_recommande} occupent surtout la zone de plausibilité dégradée.

\begin{figure}[H]
\centering
\includegraphics[width=0.78\textwidth]{fig8_stabilite_vs_physique_ffp.png}
\caption{Stabilité et plausibilité physique selon la classe FFP.}
\label{fig:stab_phys}
\end{figure}

\subsection{Agrégation au niveau station et règle du minimum}

Au niveau station, la classe retenue est la plus restrictive parmi les trois
variables. La distribution devient alors nettement plus sévère
(Tableau~\ref{tab:ffp_stations}) : seulement 10~stations sur~38 (26\,\%) restent
en usage scientifique restreint, tandis que 10~(26\,\%) sont
\texttt{non\_recommande} et 9~(24\,\%) exploratoires.

\begin{table}[H]
\centering
\caption{Classes FFP au niveau station (règle de la classe la plus restrictive).}
\label{tab:ffp_stations}
\begin{tabular}{lrr}
\toprule
Classe station & $n$ & Proportion \\
\midrule
Usage scientifique restreint & 10 & 26,3\,\% \\
Comparaison relative & 7 & 18,4\,\% \\
Suivi de tendances & 2 & 5,3\,\% \\
Exploratoire & 9 & 23,7\,\% \\
Non recommandé & 10 & 26,3\,\% \\
\bottomrule
\end{tabular}
\end{table}

Ce décalage série/station constitue un résultat central du mémoire. Il confirme
$H_3$. Une station peut afficher un score médian très élevé tout en étant
classée exploratoire, dès qu'une variable échoue. Les cas \textit{Feinstaub
Balkon}, \textit{Sportplatz Härtensdorf}, \textit{Kieseläcker} ou
\textit{emp-wetter} l'illustrent : scores médians compris entre 0,98 et 0,99,
classe station exploratoire, parce que l'humidité --- parfois la pression ---
présente une stabilité insuffisante alors que les autres variables restent
excellentes (Tableau~\ref{tab:diagnostic}).

\begin{table}[H]
\centering
\caption{Exemple de dissociation score médian élevé / classe station restrictive.}
\label{tab:diagnostic}
\begin{tabular}{llccc}
\toprule
Station & Variable & Score & Classe série & Classe station \\
\midrule
Feinstaub Balkon & humidité & 0,919 & exploratoire & exploratoire \\
Feinstaub Balkon & pression & 0,993 & usage sci. restreint & exploratoire \\
Feinstaub Balkon & température & 0,995 & usage sci. restreint & exploratoire \\
Sportplatz Härtensdorf & humidité & 0,895 & exploratoire & exploratoire \\
Sportplatz Härtensdorf & pression & 0,992 & usage sci. restreint & exploratoire \\
Sportplatz Härtensdorf & température & 0,994 & usage sci. restreint & exploratoire \\
\bottomrule
\end{tabular}
\end{table}

Les Figures~\ref{fig:carte_score} et~\ref{fig:carte_ffp} situent les stations
selon le score médian et selon la classe FFP. Elles ne prétendent pas à une
analyse spatiale formelle --- le QC spatial est désactivé --- ; elles servent à
visualiser la diversité des profils dans la zone couverte.

\begin{figure}[H]
\centering
\begin{subfigure}{0.48\textwidth}
\includegraphics[width=\textwidth]{fig5_carte_stations_confiance.png}
\caption{Score médian.}
\label{fig:carte_score}
\end{subfigure}
\hfill
\begin{subfigure}{0.48\textwidth}
\includegraphics[width=\textwidth]{fig5b_carte_stations_ffp.png}
\caption{Classe FFP station.}
\label{fig:carte_ffp}
\end{subfigure}
\caption{Répartition spatiale des stations selon la confiance et la classe d'usage.}
\end{figure}

\subsection{Réponse à $Q_3$}

La couche \textit{fit for purpose} porte une information que le score seul ne
contient pas. Elle transforme un continuum $[0,1]$ en recommandation d'usage
explicite. L'agrégation station impose ensuite une prudence systémique : pour un
usage multi-variable, la qualité de la « pire » série devient le goulot
d'étranglement. Cette propriété convient à un cadre orienté décision. Elle
reste, par construction, conservatrice --- et c'est précisément ce que $H_3$
prédisait.

%==============================================================================
\section{Discussion}
\label{sec:discussion}

\subsection{Retour sur les questions de recherche}

Sur $Q_1$, le corpus openSenseMap retenu est exploitable pour un diagnostic de
confiance : volume important, trois variables alignées, complétude médiane
élevée. Le taux de 23\,\% d'observations suspectes montre pourtant qu'un usage
naïf des brutes serait imprudent. Le cadre répond donc à un besoin réel de
filtrage et de documentation des preuves de qualité.

Sur $Q_2$, les dimensions séparent clairement les profils. L'humidité apparaît
comme la variable la plus contrainte par la stabilité ; la température offre,
dans ce corpus, le meilleur compromis entre disponibilité et confiance. Cette
différenciation justifie l'architecture multi-dimensionnelle face à un score
unique non décomposable.

Sur $Q_3$, la classification FFP et la règle du minimum au niveau station
produisent des recommandations plus restrictives que la seule médiane des
scores. Six stations conservent un score médian $\geq 0,95$ tout en restant
exploratoires ou non recommandées. Le cadre évite ainsi l'illusion d'une station
« globalement excellente » masquant une variable défaillante.

\subsection{Apports scientifiques}

Trois apports peuvent être affirmés sans surinterprétation. Premier apport : la
séparation explicite entre preuves QC (observation), score de confiance (série)
et recommandation d'usage (FFP). Deuxième apport : la reproductibilité
opérationnelle du cadre --- paramètres exportés, pipeline documenté, tables et
figures reprises ici. Troisième apport : la mise en évidence empirique d'un
goulot d'étranglement variable, principalement l'humidité via les valeurs
figées, dans un corpus PWS réel.

Ces apports concernent une confiance \textit{relative et contextuelle}. Ils ne
portent pas sur l'exactitude absolue.

\subsection{Confrontation avec la littérature}

Les résultats rejoignent plusieurs thèses établies. Teh et~al. rappellent que la
qualité des données capteurs conditionne l'utilité des applications IoT
\cite{teh2020sensor} ; le taux de 23\,\% d'observations suspectes observé ici
donne à ce constat une traduction concrète sur un corpus citoyen. De~Vos et~al.
montrent qu'un QC dédié aux PWS crowdsourcées est nécessaire avant un usage
opérationnel \cite{devos2019qc}. Notre couche de drapeaux et de statuts
s'inscrit dans cette logique, même si les variables et le barème diffèrent :
ici T/H/P et confiance relative, plutôt qu'un focus pluie. Dabrowski et~al.
insistent sur des contrôles temporels structurés en monitoring météorologique
\cite{dabrowski2022qc} ; la place accordée aux valeurs figées et aux sauts dans
notre protocole répond à cette exigence de cohérence dynamique.

Les travaux microclimatiques et prédictifs
\cite{moon2017microclimate,chithrakumar2023microclimate,yu2021lstm} confirment
l'intérêt des observations locales, tout en dépendant, en amont, de séries
exploitables. Notre contribution se situe avant la prévision : elle documente la
confiance et l'aptitude à l'usage.

Une divergence avec une partie de la littérature métrologique doit être
assumée. Sans colocalisation, nous n'avançons ni biais en °C ni RMSE. Les
études de calibration PWS--référence ne sont pas contredites ; elles
interviennent à une autre étape, en aval d'une confiance déjà documentée.

\subsection{Forces et choix assumés}

Le cadre gagne en clarté par la traçabilité des rejets, doublons et drapeaux,
ainsi que par la décomposition du score en dimensions interprétables. Les
seuils d'usage sont explicites et exportés. Le protocole tient un volume de
l'ordre de $2,4 \times 10^6$ observations de façon reproductible, ce qui dépasse
le stade du simple prototype illustratif.

Le choix de la règle du minimum au niveau station mérite d'être défendu. On
pourrait lui préférer une médiane ou une moyenne des classes. Une telle
agrégation adoucirait le verdict station, au prix d'une perte de lisibilité pour
l'utilisateur multi-variable : une humidité défaillante resterait masquée par
deux séries excellentes. Nous avons privilégié la prudence, cohérente avec
l'objectif de limiter les usages abusifs.

\subsection{Menaces à la validité}

Plusieurs menaces limitent la portée des conclusions. Sur le plan du construit,
le score mesure une confiance opérationnelle définie par le cadre, non une
vérité terrain ; le risque est de surinterpréter un score élevé comme une
exactitude. Sur le plan interne, drapeaux et seuils dépendent d'hyperparamètres
(durée minimale figée, facteur MAD, plages physiques) : d'autres réglages
modifieraient proportions suspectes et classes FFP. Sur le plan externe, le
corpus est limité à 38~stations outdoor, une zone principalement germanophone et
une période estivale d'un mois ; la généralisation à d'autres climats, saisons
ou plateformes n'est pas établie. Enfin, la corrélation de Spearman entre
complétude et score ($\rho=0,59$) décrit une association, non un lien causal ;
de même, la surreprésentation des valeurs figées en humidité n'identifie pas à
elle seule la cause matérielle.

%==============================================================================
\section{Limites de l'étude}
\label{sec:limites}

\subsection{Limites liées aux données}

Les mesures proviennent d'une plateforme citoyenne. Conditions d'installation,
entretien des capteurs, exposition réelle et interventions éventuelles de
l'hôte restent largement inconnus. Les métadonnées \texttt{lastMeasurement} de
l'API se sont révélées partiellement non fiables lors de l'acquisition ; le
protocole s'appuie donc sur les séries téléchargées, non sur ces champs.

La période d'un mois, en été, ne capture ni les régimes hivernaux ni les
épisodes extrêmes prolongés. Trente-huit stations restent modestes au regard de
l'ensemble du réseau openSenseMap.

\subsection{Limites du construit « confiance »}

Le cadre suppose que des règles de cohérence interne --- plages physiques,
continuité, stabilité, robustesse statistique --- permettent d'approcher la
confiance d'usage. L'hypothèse est raisonnable, mais incomplète. Une série peut
être internement cohérente et néanmoins biaisée : abri mal ventilé, capteur
exposé au rayonnement, dérive lente non détectée par les drapeaux retenus. La
confiance documentée ici n'élimine pas ce risque ; elle le circonscrit.

\subsection{Limites méthodologiques et métriques}

Quatre limites méthodologiques structurantes doivent être rappelées : absence de
référence colocalisée ; contrôle spatial désactivé ; seuils FFP et poids
dimensionnels choisis a priori, avec une analyse de sensibilité seulement
partielle dans le notebook~3 et non consolidée dans les exports principaux
exploités ici ; agrégation station par minimum, volontairement conservatrice.

Les taux d'anomalies et dimensions dérivées quantifient la conformité aux règles
du cadre. Ils ne remplacent pas des métriques d'erreur de mesure. Précision,
rappel ou F1 n'ont pas été calculés : aucune vérité terrain labellisée
observation par observation n'était disponible.

\subsection{Impact sur les conclusions}

Les conclusions du chapitre restent valides en tant qu'évaluation
\textbf{relative} de confiance et d'aptitude à l'usage dans le corpus étudié.
Elles ne permettent pas de classer les stations selon leur exactitude absolue,
ni d'extrapoler sans réserve à d'autres réseaux ou saisons.

%==============================================================================
\section{Perspectives}
\label{sec:perspectives}

Plusieurs prolongements découlent directement des limites précédentes. Une
validation externe, par colocalisation d'un sous-ensemble de stations avec une
référence professionnelle, permettrait d'estimer biais, RMSE et degrés de
confiance calibrés --- indicateurs absents ici pour de bonnes raisons, mais
nécessaires à l'étape suivante. Une analyse de sensibilité factorielle des poids
et seuils FFP préciserait la stabilité des classements (corrélations de rang,
taux de reclassement). Le QC spatial pourrait être réactivé lorsque la densité
locale le permet. Une extension temporelle et géographique --- année complète,
autres régions climatiques --- testerait la robustesse externe du cadre. Enfin,
l'enrichissement des métadonnées d'installation (type d'abri, hauteur,
exposition) ouvrirait la voie à des hypothèses causales plus précises sur les
valeurs figées d'humidité.

Sur le plan scientifique, le travail ouvre surtout une voie de documentation
reproductible de la confiance pour les PWS, utilisable comme étape préalable à
toute fusion de données citoyennes dans un usage analytique plus exigeant.

%==============================================================================
\section{Synthèse du chapitre}
\label{sec:synthese_ch3}

Sur un corpus de 38~stations et 2\,445\,899~observations, le cadre reproductible
proposé a permis de documenter un taux de 23,4\,\% d'observations suspectes au
QC ; d'obtenir un score de confiance médian de 0,965 au niveau des séries, avec
une dégradation spécifique de la stabilité pour l'humidité ; de classer 48\,\%
des séries en usage scientifique restreint, mais seulement 26\,\% des stations
après application de la règle du minimum ; et de montrer qu'un score médian
élevé n'implique pas une recommandation d'usage élevée.

Ces résultats répondent aux objectifs d'évaluation de la confiance dans des
données PWS low-cost, sous réserve des limites déjà discutées. Ils nourrissent
directement la conclusion générale du mémoire.
